\chapter{Stochastic Differential Equations}\label{ch:qm:stochastic-differential-equations}

At 02:00 the overnight risk run stops. A square-root variance process, stepped forward one
day at a time with the plain Euler scheme, produced a negative variance, the code took its
square root, and a NaN propagated from one path into the book's value and every limit that
depends on it. It was not one unlucky path: with the desk's calibrated parameters, three paths
in four go negative within a year, and halving the time step four times over changes nothing.
The model was not wrong, the scheme was, and the number that says so, the Feller ratio, is a
one-line function of the parameters. This chapter sets out what a \omterm{def:qm:stochastic-differential-equations:sde}{stochastic differential
equation} is and when it has a solution, the three diffusions the series uses most (geometric,
mean-reverting and square-root), the generator that links them to partial differential
equations, and the Feynman--Kac formula that turns an expectation into a PDE and back.

\section{Existence, uniqueness and the Euler scheme}

\begin{definition}[Stochastic differential equation, strong and weak solutions]\label{def:qm:stochastic-differential-equations:sde}
A \emph{stochastic differential equation}\index{stochastic differential equation} is
\[
dX_t = \mu(t, X_t)\,dt + \sigma(t, X_t)\,dW_t, \qquad X_0 = x_0,
\]
for measurable $\mu, \sigma$. A \emph{strong solution}\index{strong solution} on a given
probability space with a given \omterm{def:qm:brownian-motion:bm}{Brownian motion} $W$ is an adapted continuous process with
$X_t = x_0 + \int_0^t\mu(s, X_s)\,ds + \int_0^t\sigma(s, X_s)\,dW_s$. A \emph{weak
solution}\index{weak solution} is a pair $(X, W)$ on some filtered probability space satisfying
the equation: the \omterm{def:qm:brownian-motion:bm}{Brownian motion} is part of the answer.
\end{definition}

\begin{theorem}[Existence and uniqueness]\label{thm:qm:stochastic-differential-equations:exist}
If $|\mu(t,x) - \mu(t,y)| + |\sigma(t,x) - \sigma(t,y)| \le K|x - y|$ and $|\mu(t,x)| + |\sigma(t,x)|
\le K(1 + |x|)$, the equation has a unique \omterm{def:qm:stochastic-differential-equations:sde}{strong solution}, with $\E[\sup_{t\le T}X_t^2] <
\infty$.
\end{theorem}

\begin{proof}[Partial proof]
Picard iteration $X^{(n+1)}_t = x_0 + \int_0^t\mu(X^{(n)})\,ds + \int_0^t\sigma(X^{(n)})\,dW$: by the
isometry and Doob's inequality, $\E[\sup_{s\le t}|X^{(n+1)}_s - X^{(n)}_s|^2] \le C\int_0^t
\E[\sup_{u\le s}|X^{(n)}_u - X^{(n-1)}_u|^2]\,ds$, so the differences are bounded by $C^nt^n/n!$
and the iterates converge; Gronwall's lemma gives uniqueness (Karatzas and Shreve, 1991, §5.2).
\end{proof}

The square-root coefficient $\eta\sqrt v$ is not Lipschitz at zero, which is why the
\omterm{def:qm:stochastic-differential-equations:sqrt}{square-root process} needs its own theory below. \omterm{def:qm:stochastic-differential-equations:sde}{Weak solutions} matter because measure changes
(chapter~5) produce them: the equation $dX = \operatorname{sign}(X)\,dW$ has a \omterm{def:qm:stochastic-differential-equations:sde}{weak solution}, a
\omterm{def:qm:brownian-motion:bm}{Brownian motion}, but no strong one.

\begin{definition}[Euler--Maruyama scheme]\label{def:qm:stochastic-differential-equations:euler}
The \emph{Euler--Maruyama scheme}\index{Euler--Maruyama scheme} on the grid $t_k = k\Delta t$ is
$\hat X_{k+1} = \hat X_k + \mu(t_k, \hat X_k)\Delta t + \sigma(t_k, \hat X_k)\sqrt{\Delta t}\,Z_k$, with $Z_k$
independent standard normals.
\end{definition}

It is the discrete \omterm{def:qm:ito-calculus:integral}{Itô integral} of \cref{ch:qm:ito-calculus}: coefficients frozen at the left
point. Under the Lipschitz conditions it converges as $\Delta t \to 0$, pathwise at order
$\tfrac12$ and in law at order 1 (chapter~26 makes both orders precise). Its weakness is that
it ignores the geometry of the state space: a Gaussian increment can carry a positive process
below zero.

\section{The workhorse diffusions}

\begin{definition}[Geometric Brownian motion]\label{def:qm:stochastic-differential-equations:gbm}
A \emph{geometric Brownian motion}\index{geometric Brownian motion} solves $dS = \mu S\,dt + \sigma
S\,dW$; by \cref{ex:qm:ito-calculus:log}, $S_t = S_0\exp((\mu - \tfrac12\sigma^2)t + \sigma W_t)$, lognormal.
\end{definition}

\begin{definition}[Mean reversion, Ornstein--Uhlenbeck process, half-life]\label{def:qm:stochastic-differential-equations:ou}
A process shows \emph{mean reversion}\index{mean reversion} when its drift pulls it toward a level.
The \emph{Ornstein--Uhlenbeck process}\index{Ornstein--Uhlenbeck process} is the Gaussian case
$dX = \kappa(\bar x - X)\,dt + \sigma\,dW$, $\kappa > 0$. The \emph{half-life}\index{half-life} of a
mean-reverting process is the time $\ln 2/\kappa$ after which the expected deviation from the
level has halved.
\end{definition}

\begin{proposition}[The Ornstein--Uhlenbeck transition]\label{prop:qm:stochastic-differential-equations:ou}
$X_t = \bar x + (X_0 - \bar x)e^{-\kappa t} + \sigma\int_0^te^{-\kappa(t-s)}\,dW_s$; given $X_0$, $X_t$ is
normal with mean $\bar x + (X_0 - \bar x)e^{-\kappa t}$ and variance $\sigma^2(1 - e^{-2\kappa t})/(2\kappa)$,
tending to $\mathcal N(\bar x, \sigma^2/2\kappa)$.
\end{proposition}

\begin{proof}
Itô's formula on $e^{\kappa t}(X_t - \bar x)$ gives $d(e^{\kappa t}(X_t - \bar x)) = \sigma e^{\kappa t}\,dW_t$;
integrate, and use the isometry for the variance of the Wiener integral.
\end{proof}

The exact transition makes the \omterm{def:qm:stochastic-differential-equations:ou}{Ornstein--Uhlenbeck process} free to simulate at any step with
no error (\cref{fig:qm:stochastic-differential-equations:paths}), and its discrete sampling
is an autoregression with coefficient $e^{-\kappa\Delta t}$ (chapter~17). Spreads, funding
bases and, in One Quant Book~6, short rates are modelled this way.

\begin{definition}[Square-root process, Feller condition]\label{def:qm:stochastic-differential-equations:sqrt}
The \emph{square-root process}\index{square-root process} is $dv = \kappa(\bar v - v)\,dt + \eta\sqrt
v\,dW$ with $\kappa, \bar v, \eta > 0$ and $v_0 > 0$. The \emph{Feller condition}\index{Feller
condition} is $2\kappa\bar v \ge \eta^2$; the ratio $2\kappa\bar v/\eta^2$ is the Feller ratio.
\end{definition}

\begin{theorem}[Feller]\label{thm:qm:stochastic-differential-equations:feller}
The \omterm{def:qm:stochastic-differential-equations:sqrt}{square-root process} has a unique strong nonnegative solution. It never reaches zero if the
\omterm{def:qm:stochastic-differential-equations:sqrt}{Feller condition} holds, and reaches zero with positive probability (and is instantly reflected)
if it fails. Given $v_s$, $v_t = c\,Y$ with $Y$ noncentral chi-square with $d = 4\kappa\bar
v/\eta^2$ degrees of freedom and noncentrality $v_se^{-\kappa(t-s)}/c$, where $c = \eta^2(1 -
e^{-\kappa(t-s)})/(4\kappa)$; its stationary law is $\mathrm{Gamma}(2\kappa\bar v/\eta^2,\ \eta^2/2\kappa)$.
\end{theorem}

\admitted

Uniqueness uses the Yamada--Watanabe argument for Hölder-$\tfrac12$ coefficients; the boundary
classification is Feller's (1951), the transition law is in Cox, Ingersoll and Ross (1985).
The \omterm{def:qm:stochastic-differential-equations:sqrt}{square-root process} is the variance of the Heston model and the short rate of the CIR
model (One Quant Books 5 and 6); stochastic-volatility fits to equity smiles typically want a
high vol-of-vol $\eta$ and so violate the condition. The desk's process has $\kappa = 2$, $\bar v =
v_0 = 0.04$ and $\eta = 0.6$: a Feller ratio of 0.44, a stationary law with mean 0.04 and
standard deviation 0.06, and 27.9\% of its stationary mass below one tenth of the mean.

\begin{omfigure}
\begin{tikzpicture}
\begin{groupplot}[group style={group size=2 by 1,horizontal sep=1.7cm},
  width=7.6cm,height=4.8cm,xlabel={$t$ (years)},xmin=0,xmax=1,
  tick label style={font=\small},label style={font=\small},grid=major,grid style={gray!25},
  table/col sep=comma]
\nextgroupplot[ylabel={$X_t$},ymin=-0.35,ymax=0.6,
  yticklabel style={/pgf/number format/fixed}]
\addplot[omDef,thin] table[x=t,y=ou0]{figdata/methods/04-stochastic-differential-equations/paths.csv};
\addplot[omThm,thin] table[x=t,y=ou1]{figdata/methods/04-stochastic-differential-equations/paths.csv};
\addplot[omProp,thin] table[x=t,y=ou2]{figdata/methods/04-stochastic-differential-equations/paths.csv};
\addplot[black,thick,dashed] table[x=t,y=ou_mean]{figdata/methods/04-stochastic-differential-equations/paths.csv};
\nextgroupplot[ylabel={$v_t$},ymin=0,ymax=0.11,
  yticklabel style={/pgf/number format/fixed}]
\addplot[omDef,thin] table[x=t,y=v0]{figdata/methods/04-stochastic-differential-equations/paths.csv};
\addplot[omThm,thin] table[x=t,y=v1]{figdata/methods/04-stochastic-differential-equations/paths.csv};
\addplot[omProp,thin] table[x=t,y=v2]{figdata/methods/04-stochastic-differential-equations/paths.csv};
\addplot[black,dashed,domain=0:1] {0.04};
\end{groupplot}
\end{tikzpicture}

\omcaption{Left: three Ornstein--Uhlenbeck paths from 0.5 with $\kappa = 2$ (\omterm{def:qm:stochastic-differential-equations:ou}{half-life} 0.35 years),
$\bar x = 0$, $\sigma = 0.2$, and their expected path $0.5e^{-2t}$ (dashed). Right: three paths of the
desk's \omterm{def:qm:stochastic-differential-equations:sqrt}{square-root process} ($\kappa = 2$, $\bar v = 0.04$, $\eta = 0.6$), sampled exactly each
day; with a Feller ratio of 0.44 they spend long spells near zero. Data: the chapter's tutorial,
seeded.}\label{fig:qm:stochastic-differential-equations:paths}
\end{omfigure}

The plain Euler step from a small $v$ is Gaussian with mean $v + \kappa(\bar v - v)\Delta t$ and
standard deviation $\eta\sqrt{v\Delta t}$: from $v = 0.004$ one daily step goes negative with
probability 3.6\%, and a path spends many days near zero.
\Cref{fig:qm:stochastic-differential-equations:feller} counts the paths that produce a
negative variance within a year: 75\% at the desk's parameters, and, what surprises most, the
same 75\% with four or sixteen steps a day. Refining cannot help, because the exact process
itself comes arbitrarily close to zero on those paths; only a scheme that respects the
boundary can. Even at a Feller ratio of exactly one, a quarter of the daily Euler paths go
negative: the condition protects the process, not the scheme.

\begin{omfigure}
\begin{tikzpicture}
\begin{axis}[width=11cm,height=5cm,xmode=log,log basis x=2,
  xlabel={Feller ratio $2\kappa\bar v/\eta^2$ ($\kappa = 2$, $\bar v = 0.04$)},ylabel={paths with a negative $v$},
  tick label style={font=\small},label style={font=\small},
  xmin=0.14,xmax=4.5,ymin=0,ymax=1,grid=major,grid style={gray!25},
  xtick={0.25,0.5,1,2,4},xticklabels={0.25,0.5,1,2,4},
  legend columns=2,legend style={font=\footnotesize,at={(0.5,-0.3)},anchor=north,draw=none,fill=none,/tikz/every even column/.append style={column sep=8pt}},
  table/col sep=comma]
\addplot[omThm,very thick,mark=*,mark size=1.6pt] table[x=ratio,y=daily]{figdata/methods/04-stochastic-differential-equations/feller.csv};
\addplot[omDef,thick,dashed,mark=triangle*,mark size=1.8pt] table[x=ratio,y=fine]{figdata/methods/04-stochastic-differential-equations/feller.csv};
\addplot[gray,thin] coordinates {(1,0) (1,1)};
\addplot[black,only marks,mark=o,mark size=3pt] coordinates {(0.4444,0.752)};
\node[font=\footnotesize,anchor=west] at (axis cs:0.5,0.93) {desk: 0.44, 75\%};
\legend{daily steps,four steps a day}
\end{axis}
\end{tikzpicture}

\omcaption{Share of plain-Euler paths of the \omterm{def:qm:stochastic-differential-equations:sqrt}{square-root process} that produce a negative
variance within one year, against the Feller ratio (vol-of-vol $\eta$ from 0.2 to 1.0). Refining the
step does not remove the failures; exact and full-truncation schemes have none. Data: 20\,000
and 10\,000 seeded paths per point.}\label{fig:qm:stochastic-differential-equations:feller}
\end{omfigure}

Two fixes are standard. The exact scheme samples the noncentral chi-square transition of
\cref{thm:qm:stochastic-differential-equations:feller} and is positive by construction.
\emph{Full truncation} (Lord, Koekkoek and van Dijk, 2010) keeps the Euler step but uses $v^+ =
\max(v, 0)$ in both coefficients, reporting $v^+$: it never takes the square root of a negative
number, and its bias vanishes as $\Delta t \to 0$. One Quant Book~5, chapter~23, adds the
quadratic-exponential scheme used for Heston in production.

\section{Generators and the Kolmogorov equations}

\begin{definition}[Infinitesimal generator]\label{def:qm:stochastic-differential-equations:generator}
The \emph{infinitesimal generator}\index{infinitesimal generator} of a time-homogeneous
diffusion $dX = \mu(X)\,dt + \sigma(X)\,dW$ is the operator $\mathcal Lf(x) = \lim_{t\downarrow 0}
(\E_x[f(X_t)] - f(x))/t = \mu(x)f'(x) + \tfrac12\sigma^2(x)f^{\prime\prime}(x)$ on $C^2$ functions.
\end{definition}

\begin{proposition}[Dynkin's formula]\label{prop:qm:stochastic-differential-equations:dynkin}
For $f \in C^2$ with compact support and a \omterm{def:qm:probability-at-speed:stopping}{stopping time} $\tau$ with $\E_x[\tau] < \infty$,
$\E_x[f(X_\tau)] = f(x) + \E_x\int_0^\tau\mathcal Lf(X_s)\,ds$.
\end{proposition}

\begin{proof}
By Itô's formula, $f(X_t) - f(x) - \int_0^t\mathcal Lf(X_s)\,ds = \int_0^tf'(X_s)\sigma(X_s)\,dW_s$, a
\omterm{def:qm:probability-at-speed:martingale}{martingale} with bounded integrand; apply optional stopping at $\tau \wedge n$ and let $n \to \infty$.
\end{proof}

\begin{definition}[Kolmogorov equations, stationary distribution]\label{def:qm:stochastic-differential-equations:kolmogorov}
For $u(t, x) = \E[g(X_T) \mid X_t = x]$, the \emph{Kolmogorov backward equation}\index{Kolmogorov
backward equation} is $\partial_tu + \mathcal Lu = 0$ with $u(T, \cdot) = g$. The transition density
$p(t, y)$ of $X_t$ solves the \emph{Kolmogorov forward equation}\index{Kolmogorov forward equation},
also called the \emph{Fokker--Planck equation}\index{Fokker--Planck equation}, $\partial_tp =
\mathcal L^*p = -\partial_y(\mu p) + \tfrac12\partial_{yy}(\sigma^2p)$. A \emph{stationary
distribution}\index{stationary distribution} of a \omterm{def:qm:brownian-motion:markov}{Markov process} is a law that, taken as the law of
$X_0$, is the law of every $X_t$; for a diffusion its density solves $\mathcal L^*p = 0$.
\end{definition}

The backward equation looks at a payoff from where one stands; the forward equation pushes a
density forward in time. For the \omterm{def:qm:stochastic-differential-equations:ou}{Ornstein--Uhlenbeck process}, $\mathcal L^*p = 0$ integrates once to
$\kappa(\bar x - y)p = \tfrac12\sigma^2p'$, whose solution is the $\mathcal N(\bar x, \sigma^2/2\kappa)$ density; for the
\omterm{def:qm:stochastic-differential-equations:sqrt}{square-root process}, $\kappa(\bar v - y)p = \tfrac12\eta^2(yp)'$ gives $p \propto y^{2\kappa\bar
v/\eta^2 - 1}e^{-2\kappa y/\eta^2}$, the Gamma law of \cref{thm:qm:stochastic-differential-equations:feller}.
When the Feller ratio is below one the exponent is negative and the density is infinite at
zero (\cref{fig:qm:stochastic-differential-equations:stationary}).

\begin{omfigure}
\begin{tikzpicture}
\begin{groupplot}[group style={group size=2 by 1,horizontal sep=1.7cm},
  width=7.6cm,height=4.8cm,xlabel={$v$},xmin=0,xmax=0.16,ymin=0,
  tick label style={font=\small},label style={font=\small},grid=major,grid style={gray!25},
  xticklabel style={/pgf/number format/fixed},scaled x ticks=false,
  table/col sep=comma]
\nextgroupplot[ylabel={density},ymax=66,title={\small desk: Feller ratio 0.44}]
\addplot[omDef,thick,const plot mark mid] table[x=v,y=desk_hist]{figdata/methods/04-stochastic-differential-equations/stationary.csv};
\addplot[omThm,only marks,mark=*,mark size=1.3pt] table[x=v,y=desk_pdf]{figdata/methods/04-stochastic-differential-equations/stationary.csv};
\nextgroupplot[ymax=26,title={\small $\eta = 0.2$: Feller ratio 4}]
\addplot[omDef,thick,const plot mark mid] table[x=v,y=f2_hist]{figdata/methods/04-stochastic-differential-equations/stationary.csv};
\addplot[omThm,only marks,mark=*,mark size=1.3pt] table[x=v,y=f2_pdf]{figdata/methods/04-stochastic-differential-equations/stationary.csv};
\end{groupplot}
\end{tikzpicture}

\omcaption{Stationary laws of the \omterm{def:qm:stochastic-differential-equations:sqrt}{square-root process} with $\kappa = 2$, $\bar v = 0.04$: histograms of
50\,000 exact paths after five years (steps) against the Gamma law that solves $\mathcal L^*p = 0$,
averaged over each bin (dots). With the desk's $\eta = 0.6$ the density piles up at zero; with
$\eta = 0.2$ zero is unattainable. Data: the chapter's tutorial, seeded.}\label{fig:qm:stochastic-differential-equations:stationary}
\end{omfigure}

\section{Feynman--Kac}

\begin{theorem}[Feynman--Kac]\label{thm:qm:stochastic-differential-equations:fk}
Let $X$ solve the equation of \cref{def:qm:stochastic-differential-equations:sde}, $r \ge 0$ and $g$
continuous and of polynomial growth, and let $u \in C^{1,2}$ of polynomial growth solve
\[
\partial_tu + \mu\,\partial_xu + \tfrac12\sigma^2\partial_{xx}u - r(t,x)u = 0, \qquad u(T, x) = g(x).
\]
Then $u(t, x) = \E\bigl[e^{-\int_t^Tr(s, X_s)\,ds}g(X_T) \bigm| X_t = x\bigr]$.
\end{theorem}

\begin{proof}
With $D_s = e^{-\int_t^sr(u, X_u)\,du}$, Itô's product rule gives $d(D_su(s, X_s)) = D_s(\partial_su +
\mathcal Lu - ru)\,ds + D_s\sigma\partial_xu\,dW_s = D_s\sigma\partial_xu\,dW_s$. The growth conditions make
the stochastic integral a \omterm{def:qm:probability-at-speed:martingale}{martingale}, so $u(t, x) = \E[D_Tu(T, X_T) \mid X_t = x]$.
\end{proof}

The theorem runs both ways: an expectation can be computed by solving a PDE (chapter~27), and a
PDE by simulating paths (chapter~26). \Cref{fig:qm:stochastic-differential-equations:map}
collects the links. With $X = r$ itself an Ornstein--Uhlenbeck short rate and $g = 1$, $u$ is the
price of a zero-coupon bond, and the PDE has the affine solution $u = e^{A(\tau) - B(\tau)r}$ with
$B(\tau) = (1 - e^{-\kappa\tau})/\kappa$ and $A(\tau) = (\bar r - \sigma^2/2\kappa^2)(B - \tau) - \sigma^2B^2/4\kappa$,
found by substituting and matching the terms in $r$. With $r_0 = 3\%$, $\bar r = 4\%$, $\kappa = 0.5$
and $\sigma = 1\%$, the five-year price is 0.8343, and 20\,000 simulated paths give $0.8341 \pm 0.0002$.
One Quant Book~6, chapter~7, turns this calculation into a model of the curve.

\begin{omfigure}
\begin{tikzpicture}[font=\small,node distance=0.7cm and 2.7cm,
  bx/.style={draw=omDef,rounded corners=2pt,fill=omDef!6,inner sep=4pt,align=center,text width=4.2cm}]
\node[bx] (sde) {SDE\\$dX = \mu\,dt + \sigma\,dW$};
\node[bx,right=of sde] (gen) {generator\\$\mathcal Lf = \mu f' + \tfrac12\sigma^2f^{\prime\prime}$};
\node[bx,below=of gen] (bwd) {backward equation\\$\partial_tu + \mathcal Lu - ru = 0$, $u(T) = g$};
\node[bx,below=of sde] (exp) {expectation\\$u = \E_{t,x}[e^{-\int r}g(X_T)]$};
\node[bx,below=of bwd] (fwd) {forward equation\\$\partial_tp = \mathcal L^*p$\\stationary: $\mathcal L^*p = 0$};
\node[bx,below=of exp] (law) {law of $X_t$\\density $p(t, y)$};
\draw[-{Stealth},thick] (sde) -- node[above,font=\footnotesize] {Itô} (gen);
\draw[-{Stealth},thick] (gen) -- (bwd);
\draw[{Stealth}-{Stealth},thick,omThm] (bwd) -- node[above,font=\footnotesize,omThm] {Feynman--Kac} (exp);
\draw[-{Stealth},thick] (gen.east) to[bend left=35] node[right,font=\footnotesize] {adjoint} (fwd.east);
\draw[{Stealth}-{Stealth},thick] (fwd) -- (law);
\draw[-{Stealth},thick,gray] (sde) -- node[left,font=\footnotesize] {simulate} (exp);
\end{tikzpicture}

\omcaption{How a diffusion's equations fit together. The generator comes from Itô's formula;
expectations of payoffs solve the backward equation (Feynman--Kac, with discounting), densities
solve the forward equation, and stationary laws solve $\mathcal L^*p = 0$.}\label{fig:qm:stochastic-differential-equations:map}
\end{omfigure}

\section{Tutorial: the overnight NaN}\label{tut:qm:stochastic-differential-equations:nan}

\begin{tutorial}
\textbf{Goal.} Reproduce the overnight failure, measure it against the Feller ratio, fix it two
ways, and check the stationary law and a Feynman--Kac price. \textbf{End state:}
\cref{fig:qm:stochastic-differential-equations:paths,fig:qm:stochastic-differential-equations:feller,fig:qm:stochastic-differential-equations:stationary}
and the numbers of the weekend problem.
\begin{tutsteps}
\item \textbf{The schemes.} The running project samples the \omterm{def:qm:stochastic-differential-equations:sqrt}{square-root process} exactly and by
Euler, plain or with full truncation.
\omcode{code/firm/mcengine/firm_mcengine.py}{150}{180}{Exact and Euler steps of the square-root process.}\label{lst:qm:stochastic-differential-equations:sqrt}
\item \textbf{The count}: the share of plain-Euler paths that produce a negative variance, and a
Feynman--Kac check by simulation.
\omcode{code/methods/04-stochastic-differential-equations/python/qm_sde.py}{77}{90}{A zero-coupon price under an Ornstein--Uhlenbeck short rate, by simulation and by the PDE.}\label{lst:qm:stochastic-differential-equations:fk}
\item \textbf{Run} \texttt{problem()}, \texttt{feller\_table()}, \texttt{stationary\_histograms()} and
\texttt{fig\_sde.py}; the C++20 and Rust twins of the Euler and Ornstein--Uhlenbeck steps are in
\texttt{code/firm/mcengine/}.
\end{tutsteps}
\textbf{What to change next.} Add the reflection scheme ($v \leftarrow |v|$) and compare its one-year
mean with the exact scheme's; give the variance process a time-dependent level $\bar v(t)$ and
check the generator's prediction for $\E[v_t]$.
\end{tutorial}

\section{Build: the Monte Carlo engine, stage two}\label{bld:qm:stochastic-differential-equations:mcengine}

\begin{build}
\bfield{Purpose} Step the diffusions every later chapter simulates, with the scheme each one
needs, and never return a NaN.
\bfield{Interface} \texttt{euler\_maruyama(mu, sigma, x0, T, n\_steps, n\_paths, seed)};
\texttt{ou\_exact(x0, kappa, xbar, sigma, T, n\_steps, n\_paths, seed)}; \texttt{sqrt\_exact(v0,
kappa, vbar, eta, \dots)}; \texttt{sqrt\_euler(\dots, scheme)} with \texttt{plain} or
\texttt{full\_truncation}; \texttt{feller\_ratio(kappa, vbar, eta)}. C++20 and Rust:
\texttt{ou\_exact\_path} and \texttt{sqrt\_euler\_path} over the stage-one \texttt{NormalStream}.
\bfield{Rules} Parameters in the series notation ($\kappa$, $\bar x$, $\bar v$, $\eta$); exact transitions where
they exist; a scheme that can leave the state space is named as such and never the default.
\bfield{Acceptance tests} Stationary mean and variance of both processes; Euler converges to the
exact Ornstein--Uhlenbeck mean; the exact square-root scheme is nonnegative with the right
moments; plain Euler fails at the desk's parameters and full truncation never does, in all three
languages.
\bfield{Stretch} The quadratic-exponential scheme; a Milstein step (chapter~26).
\end{build}

\begin{omsources}
\item W.~Feller, ``Two singular diffusion problems'', \emph{Annals of Mathematics} 54, 1951.
\item J.~C.~Cox, J.~E.~Ingersoll and S.~A.~Ross, ``A theory of the term structure of interest
rates'', \emph{Econometrica} 53, 1985.
\item G.~E.~Uhlenbeck and L.~S.~Ornstein, ``On the theory of the Brownian motion'',
\emph{Physical Review} 36, 1930.
\item M.~Kac, ``On distributions of certain Wiener functionals'', \emph{Transactions of the AMS}
65, 1949.
\item R.~Lord, R.~Koekkoek and D.~van Dijk, ``A comparison of biased simulation schemes for
stochastic volatility models'', \emph{Quantitative Finance} 10, 2010.
\end{omsources}

\section{Exercises}

\begin{exercise}[$\star$]\label{exo:qm:stochastic-differential-equations:1}
A spread follows an \omterm{def:qm:stochastic-differential-equations:ou}{Ornstein--Uhlenbeck process} with $\kappa = 2$ per year. What is its \omterm{def:qm:stochastic-differential-equations:ou}{half-life}
in years and in trading days?
\end{exercise}

\begin{exercise}[$\star$]\label{exo:qm:stochastic-differential-equations:2}
A stock follows a \omterm{def:qm:stochastic-differential-equations:gbm}{geometric Brownian motion} with $\mu = 5\%$ and $\sigma = 20\%$. What is the probability
that it is below its starting price after one year?
\end{exercise}

\begin{exercise}[$\star$]\label{exo:qm:stochastic-differential-equations:3}
Write the generator of the \omterm{def:qm:stochastic-differential-equations:ou}{Ornstein--Uhlenbeck process}, apply it to $f(x) = x$, and deduce the ODE
satisfied by $\E[X_t]$.
\end{exercise}

\begin{exercise}[$\star\star$]\label{exo:qm:stochastic-differential-equations:4}
From the forward equation, find the stationary variance of the \omterm{def:qm:stochastic-differential-equations:ou}{Ornstein--Uhlenbeck process} with
$\kappa = 2$, $\sigma = 0.4$.
\end{exercise}

\begin{exercise}[$\star\star$]\label{exo:qm:stochastic-differential-equations:5}
With Dynkin's formula, compute $\E[v_{0.5}]$ for the \omterm{def:qm:stochastic-differential-equations:sqrt}{square-root process} with $v_0 = 0.09$, $\bar v =
0.04$, $\kappa = 2$.
\end{exercise}

\begin{exercise}[$\star\star$]\label{exo:qm:stochastic-differential-equations:6}
Check the five-year zero-coupon price of the chapter (0.8343) from the formula for $A$ and $B$.
\end{exercise}

\begin{exercise}[$\star\star\star$]\label{exo:qm:stochastic-differential-equations:7}
\emph{Coding.} With \texttt{negative\_fraction}, find the share of daily plain-Euler paths that go
negative within a year when the Feller ratio is exactly one ($\eta = 0.4$), and the largest $\eta$ in
the table for which it is below 1\%.
\end{exercise}

\begin{exercise}[$\star\star\star$]\label{exo:qm:stochastic-differential-equations:8}
\emph{Find the flaw.} ``We floor the variance at zero after every Euler step, so the scheme is
now exact.''
\end{exercise}

\section{Problem: The NaN in the Overnight Run}

\begin{problem}[{Weekend problem --- a variance process that violates the Feller condition}]\label{pb:qm:stochastic-differential-equations:1}
The desk's variance process is $dv = \kappa(\bar v - v)\,dt + \eta\sqrt v\,dW$ with $v_0 = \bar v = 0.04$,
$\kappa = 2$ and $\eta = 0.6$, simulated for one year of 252 days.

\medskip\noindent\textbf{Part I --- The process.}
\begin{enumerate}
\item What is the Feller ratio? Is zero attainable?
\item What are the stationary law, its mean and standard deviation?
\item What share of the stationary mass lies below 0.004?
\item What vol-of-vol would make the Feller ratio exactly one?
\item What is the \omterm{def:qm:stochastic-differential-equations:ou}{half-life} of the expected variance's return to $\bar v$?
\end{enumerate}

\medskip\noindent\textbf{Part II --- Plain Euler.}
\begin{enumerate}[resume]
\item From $v = 0.004$, what is the probability that one daily Euler step goes negative?
\item What share of paths produce a negative variance within the year?
\item With four and sixteen steps a day, and with weekly steps?
\item Why does refining the step not help?
\item What does one NaN do to a risk run that sums over paths?
\end{enumerate}

\medskip\noindent\textbf{Part III --- Fixes.}
\begin{enumerate}[resume]
\item What minimum and one-year mean does the exact scheme give?
\item What one-year mean does full truncation give, and on what share of steps is it at zero?
\item Which of the two is biased, and how does the bias behave as $\Delta t \to 0$?
\item What would reflection ($v \leftarrow |v|$) do to the mean near zero?
\item Why do production Heston engines use another scheme?
\end{enumerate}

\medskip\noindent\textbf{Part IV --- Judgement.}
\begin{enumerate}[resume]
\item Should the desk impose the \omterm{def:qm:stochastic-differential-equations:sqrt}{Feller condition} in its calibration?
\item What should an automated test of the simulation engine assert?
\item What belongs in the model-review note about this incident?
\item State the \emph{named result}: the Feller ratio and the share of failing paths.
\item In one sentence: what does the \omterm{def:qm:stochastic-differential-equations:sqrt}{Feller condition} protect, and what does it not?
\end{enumerate}
\end{problem}

\section{Interview questions}

\begin{interviewq}[$\star$ \iqroles{trader, researcher}]\label{iq:qm:stochastic-differential-equations:1}
Solve $dS = \mu S\,dt + \sigma S\,dW$. What is the median of $S_T$?
\end{interviewq}

\begin{interviewq}[$\star$ \iqroles{researcher}]\label{iq:qm:stochastic-differential-equations:2}
A spread mean-reverts with a \omterm{def:qm:stochastic-differential-equations:ou}{half-life} of 10 days. What is $\kappa$, and how much of today's
deviation is expected to remain after 30 days?
\end{interviewq}

\begin{interviewq}[$\star\star$ \iqroles{researcher, bank}]\label{iq:qm:stochastic-differential-equations:3}
What is the \omterm{def:qm:stochastic-differential-equations:sqrt}{Feller condition}, and what happens to a \omterm{def:qm:stochastic-differential-equations:sqrt}{square-root process} and to its simulation
when it fails?
\end{interviewq}

\begin{interviewq}[$\star\star$ \iqroles{researcher, bank}]\label{iq:qm:stochastic-differential-equations:4}
State the Feynman--Kac formula and use it to write the PDE for $\E[e^{-rT}g(S_T)]$ under a
\omterm{def:qm:stochastic-differential-equations:gbm}{geometric Brownian motion}.
\end{interviewq}

\begin{interviewq}[$\star\star$ \iqroles{developer}]\label{iq:qm:stochastic-differential-equations:5}
The overnight Monte Carlo sometimes returns NaN. How do you find the cause and make it
impossible?
\end{interviewq}

\begin{interviewq}[$\star\star\star$ \iqroles{researcher}]\label{iq:qm:stochastic-differential-equations:6}
What is the difference between a strong and a \omterm{def:qm:stochastic-differential-equations:sde}{weak solution}? Give an equation with one but not
the other.
\end{interviewq}
